\newcommand*{\path7}{chapters/chapter7}

\section{Topographical Descriptors}

Physicochemical and topological descriptors have limitations since they don't take into account the three-dimensional structure of molecules.
Binding sites exists in three dimensions, therefore the shape complementarity between the binding site and a molecule is important in order to create molecular interactions. 
Topographical (or molecular shape) descriptors should account for the interplay between molecular interactions and three dimensional structure.

The similarity principle also apply to the three-dimensional shape of molecules.
The fact that molecules can change shape over time and assume different conformations poses complications on the use of topographical descriptors and is known as the conformer problem.


%%%%%%%%%%%%%%%%%%%%%%%%%%%%%%%%%%%%%%
\subsection{Topographical Descriptors}

Geometric descriptor can be computed from atomic coordinates.
There are thw types of geometric descriptors:
\begin{itemize}
    \item Cartesian descriptor (external, x3D),
    \item Internal descriptor (i3D).
\end{itemize}
Cartesian descriptors depend on a reference frame and require some sort of alignment or normalization to be used; they are alignment-dependent. 
Internal descriptors are based on relative distances between atoms; they are alignment-free.

The principal moments of inertia (PMI) descriptor consists in the first three principal moments of inertia of the molecule and contain information about the general shape of the molecule (rod-like, disc-like or sphere-like). 
One limitation of PMI is that is completely size-independent.

For the plane of best fir (PBF) descriptor, the heavy atoms of a molecule are considere data points and the plane of best fit is computed using a least-square fit. 
The PBF descriptor is the average error of the fit (i.e. the average distance of heavy atoms from the plane).

There is only a weak correlation between PMI and PBF. The PBF descriptor is much more discriminatory for planar or near-planar molecules than PMI.


%%%%%%%%%%%%%%%%%%%%%%%%%%%
\subsection{Pharmacophores}

A pharmacophore is a collection of steric and electronic features necessary for optimal supramolecular interactions with a specific biological target and to trigger its biological response.

To generate a pharmacophore, one start with one active ligand or a set of active ligands (and eventually some inactive ones) and proceeds as follows:
\begin{enumerate}
    \item Enumerate lowe energy the conformers for each ligand,
    \item Align all the conformers of all the ligand using a superimposition algorithm,
    \item Select the best overlay with a single conformer of each active ligand as optimal overlay,
    \item Abstract the optimal overlay to get the pharmacophore model (reduce atomic and structural features to feature points and spheres).
\end{enumerate}

Once a pharmacophore is generated, it needs to be validated before being applied in virtual screening.


%%%%%%%%%%%%%%%%%%%%%%%%%%%%%%%%%%%%%%%%%%%%%%%%%%%%%%%
\subsection{ROCS: Rapid Overlay of Chemical Structures}

In the rapid overlay of chemical structures (ROCS) method, Gaussian functions represent the space occupied by atoms and different colours are given to different chemical types (rings, HBA, HBD, \dots).


%%%%%%%%%%%%%%%%%%%%%%%%%%%%%%%%%%%%%%%%%%%%%
\subsection{USR: Ultrafast Shape Recognition}

The ultrafast shape recognition is a very fast method to compute the similarity between three dimensional structures.
In this method, the centroid atom and some other atoms are used to compute statistical moments about atomic distribution to encode the three dimensional shape of the molecule.
The similarity between two molecules is finally quantified as the Manhattan distance between the vectors of statistical moments.


%%%%%%%%%%%%%%%%%%%%%%%%%%%%%%%%%%%%%%%%%%%%%%%%
\subsection{XED: Extended Electron Distribution}

The extended electron distribution (XED) model provides a more accurate description of the electron distribution around atoms, by capturing more complex features such as lone pairs.


%%%%%%%%%%%%%%%%%%%%%%%%%%%%%%%%%%%%%%%%%%%%%%%%%%%%%%%%%%%
\subsection{Conformer Generation and the Conformer Problem}

Conformer generation is a challenging and time-consuming process. 
There are very many conformers for a given molecule, in addition to tautomers and stereoisomers.
There are two main ways of generating conformers:
\begin{itemize}
    \item Systematic search,
    \item Stochastic (random) search.
\end{itemize}

The conformer problems comes from the fact that many conformers can be generated for a given molecule and is not clear which one is the most appropriate for drug design.

Systematic search consider each rotatable bond in the molecule and explore the torsional angles of each bond.
The problem with systematic search is the combinatorial explosion of all possible combinations of discrete bond rotations.
Many conformers generated by systematic search will introduce steric clashes or high-energy conformations.
Some methods are able to avoid the generation of such conformations, thus lowering the number of conformers to be generated. Other methods treat the molecule as a set of rigid fragments and only the rotatable bonds between fragments are considered.

Stochastic search randomly explore the space of rotation of the rotatable bonds. 
Since the search is heuristic, is not clear when it should be terminated.

Once conformers are generated, it is difficult or impossible to know if any structure is similar to the one expected during a binding event. This is known as the conformer problem.